% Options for packages loaded elsewhere
\PassOptionsToPackage{unicode}{hyperref}
\PassOptionsToPackage{hyphens}{url}
\PassOptionsToPackage{dvipsnames,svgnames,x11names}{xcolor}
\PassOptionsToPackage{space}{xeCJK}
\documentclass[
]{article}
\usepackage{xcolor}
\usepackage[margin=2.5cm]{geometry}
\usepackage{amsmath,amssymb}
\setcounter{secnumdepth}{-\maxdimen} % remove section numbering
\usepackage{iftex}
\ifPDFTeX
  \usepackage[T1]{fontenc}
  \usepackage[utf8]{inputenc}
  \usepackage{textcomp} % provide euro and other symbols
\else % if luatex or xetex
  \usepackage{unicode-math} % this also loads fontspec
  \defaultfontfeatures{Scale=MatchLowercase}
  \defaultfontfeatures[\rmfamily]{Ligatures=TeX,Scale=1}
\fi
\usepackage{lmodern}
\ifPDFTeX\else
  % xetex/luatex font selection
    \setmainfont[]{DejaVu Sans}
    \setmonofont[]{DejaVu Sans Mono}
  \ifXeTeX
    \usepackage{xeCJK}
    \setCJKmainfont[]{Noto Sans CJK SC}
          \fi
  \ifLuaTeX
    \usepackage[]{luatexja-fontspec}
    \setmainjfont[]{Noto Sans CJK SC}
  \fi
\fi
% Use upquote if available, for straight quotes in verbatim environments
\IfFileExists{upquote.sty}{\usepackage{upquote}}{}
\IfFileExists{microtype.sty}{% use microtype if available
  \usepackage[]{microtype}
  \UseMicrotypeSet[protrusion]{basicmath} % disable protrusion for tt fonts
}{}
\makeatletter
\@ifundefined{KOMAClassName}{% if non-KOMA class
  \IfFileExists{parskip.sty}{%
    \usepackage{parskip}
  }{% else
    \setlength{\parindent}{0pt}
    \setlength{\parskip}{6pt plus 2pt minus 1pt}}
}{% if KOMA class
  \KOMAoptions{parskip=half}}
\makeatother
\usepackage{longtable,booktabs,array}
\usepackage{calc} % for calculating minipage widths
% Correct order of tables after \paragraph or \subparagraph
\usepackage{etoolbox}
\makeatletter
\patchcmd\longtable{\par}{\if@noskipsec\mbox{}\fi\par}{}{}
\makeatother
% Allow footnotes in longtable head/foot
\IfFileExists{footnotehyper.sty}{\usepackage{footnotehyper}}{\usepackage{footnote}}
\makesavenoteenv{longtable}
\setlength{\emergencystretch}{3em} % prevent overfull lines
\providecommand{\tightlist}{%
  \setlength{\itemsep}{0pt}\setlength{\parskip}{0pt}}
\usepackage{bookmark}
\IfFileExists{xurl.sty}{\usepackage{xurl}}{} % add URL line breaks if available
\urlstyle{same}
\hypersetup{
  colorlinks=true,
  linkcolor={Maroon},
  filecolor={Maroon},
  citecolor={Blue},
  urlcolor={Blue},
  pdfcreator={LaTeX via pandoc}}

\author{}
\date{}

\begin{document}

\section{统一框架中的神经宏编译：把慢构造路径编译为快推理路径}\label{ux7edfux4e00ux6846ux67b6ux4e2dux7684ux795eux7ecfux5b8fux7f16ux8bd1ux628aux6162ux6784ux9020ux8defux5f84ux7f16ux8bd1ux4e3aux5febux63a8ux7406ux8defux5f84}

\section{Neural Macro Compilation in the Unified Framework of
Representation, Logic, and Intuition --- Intuition Searches
Form}\label{neural-macro-compilation-in-the-unified-framework-of-representation-logic-and-intuition-intuition-searches-form}

\begin{quote}
\textbf{中文版 v0.1} \textbar{} 2026-08-11 (翻译自英文版) 目标:
ICML/COLM (方法论文) \textbar{} 专利: M1 (语义路径编译)
\textbf{框架定位}: 属于《表示、逻辑、直觉的统一框架》(纲领:
\texttt{Unified\_Framework\_Representation\_Logic\_Intuition.md})。本文确立
\textbf{构造→直觉} 的方法环节:
在同一模型内把慢构造路径编译为快推理路径。姊妹篇: P1 (形式→构造理论), P0
(三通道等价认知)。 证据: exp50/51/80/90, 经 run\_exp 官方判定
(docs/paper\_data/)
\end{quote}

\begin{center}\rule{0.5\linewidth}{0.5pt}\end{center}

\textbf{Repository:
\href{https://github.com/YuchenWang-ai/Neural-Macro-Compilation-Compiling-Slow-Construction-Paths-into-Fast-Reasoning-Paths}{YuchenWang-ai/Neural-Macro-Compilation-Compiling-Slow-Construction-Paths-into-Fast-Reasoning-Paths}}

\subsection{摘要}\label{ux6458ux8981}

一个已学会基础运算的模型, 往往\emph{已经能够执行}新的复合算子 ---
通过展开其定义链 --- 但沿着一条长构造路径
(许多步骤/token)。我们提出\emph{神经宏编译}: 在\emph{同一模型}内,
用少量自动生成的样本, 把慢构造路径编译为快推理路径, 同时 (i)
精确保持语义 (可验证, 带回退) 且 (ii) 绝不扩展语义闭包。在一个 2
层、\textless1 MB 的 transformer 上: 2000 位进位链在构造 (0.62 s)、形式
(0.069 s)、直觉 (0.086 s) 三个通道上执行结果一致 (1.000); 编译的快路径对
2000 位、60 进制、100\% 数字匿名化零衰减泛化; 跳步可见为 20 位加法的
82→5 个 token。与蒸馏不同, 教师和学生是同一模型的慢路径与快路径
(构造通道自编译); 与程序合成不同, 没有程序被带外编译。在统一框架内,
本文是 \textbf{构造→直觉} 的方法环节: 精确的逻辑构造被编译为闪念直觉,
再由直觉搜索形式 --- 闭合形式 → 构造 → 直觉 → 形式的循环。

\subsection{1 引言}\label{ux5f15ux8a00}

\textbf{问题。} 当模型学会基础运算后, 新算子 σ 已经可执行:
其定义链展开为模型能跟随的基础运算。但路径很长 --- 逐位迭代、长 token
序列 --- 因此推理成本随构造长度而非语义规模增长。

\textbf{现有补救是侵入式的。}
蒸馏需要独立的教师和新参数。少样本微调有破坏已学能力的风险且无等价保证。程序到神经编译
(ICML 2024)
把程序当作外部输入并带外编译。它们都没有利用模型\emph{已经拥有}语义这一事实。

\textbf{本工作。} 同一模型已经包含 σ 的两条路径: 慢的\emph{构造}路径
(定义驱动, 语义完备)
和一条训练可在权重中建立的快路径。我们把前者编译为后者 ---
\emph{神经宏编译} --- 监督信号由慢路径自身自动生成,
配验证器与运行时回退。宏增加的是速度, 不是语义。与 P1 的关系: P1
确立表示\emph{支持}系统化泛化\emph{及其原因} (形式化的归纳);
本文确立一个已可执行的算子\emph{如何}被重新编译为廉价路径 --- 两者互补,
P2 不重新推导 P1 的泛化结果。

\textbf{贡献。} - \textbf{C1 --- 语义路径自编译。}
用少量自生成样本把模型自身的慢构造路径编译为快推理路径,
保持语义与语义闭包。 - \textbf{C2 --- 零样本语义闭包。}
新算子零训练即可执行 (构造路径); 训练只降低成本。 - \textbf{C3 ---
可验证等价 + 回退。} 快输出与慢输出核对; 失败路由回慢路径。 - \textbf{C4
--- 实证画像。} 2000 位进位链所有通道正确; 快路径在宽度 2000 时便宜 7×;
零衰减泛化; 82→5 token 跳步。

\subsection{2 设定}\label{ux8bbeux5b9a}

本文直接构建于 P1 的 token 原生表示之上: 概念携带结构化定义; σ
的\emph{定义链}是 σ 定义引用概念的传递闭包 (P1
§2)。编译目标是该表示之上的通道。

\textbf{慢路径 (构造通道)。} 定义驱动展开: σ
的定义链展开为逐步执行的基础运算 (逐位迭代、进位传播)。语义完备 ---
语义由定义保证; 成本随构造长度增长。

\textbf{快路径 (直觉通道)。} 编译进权重的直接执行: 短序列内输入 → 输出,
跳过中间构造步骤。成本随输出长度而非构造深度增长。

\textbf{形式路径 (重写通道)。} 定义驱动的重写引擎 --- 与构造共享定义源;
在验证中作为独立交叉核对。

\textbf{目标。} 对带展开路径 T\_σ 的新算子 σ: - Sem\_fast(σ) =
Sem\_expanded(T\_σ) --- 精确语义等价, 可验证; - Cost\_fast ≪
Cost\_expanded --- 步骤/token/延迟降低; - Closure\_fast(σ) =
Closure\_expanded(σ) --- 无语义闭包扩展。

\textbf{术语。} 快路径即\emph{推理快路径} --- 直觉: 难以获得
(编译需要充分训练; 配置相关轮数), 但易于通过跳步泛化。完整论述在 P1 (§3,
§7); 本文聚焦编译方法。

\subsection{3 方法:
神经宏编译}\label{ux65b9ux6cd5-ux795eux7ecfux5b8fux7f16ux8bd1}

五个阶段。

\textbf{S1 --- 路径识别。} 把 σ 的定义链展开为基础运算 (E(σ) = T\_σ)。若
σ 可由模型的基础运算表达, 则它可编译 --- 由定义图判定, 而非独立合成器。

\textbf{S2 --- 样本生成。} 用慢路径 (构造通道 / 参考求值器) 自动生成 (σ,
x\_i, y\_i) 对。自监督标签: 样本即模型自身的语义。少量样本即可 (数十个),
因为模型已知语义;
样本只教直接路径的\emph{形状}。编译质量对样本数的依赖是可证伪的主张, 在
§5 中检验 (当前 pending)。

\textbf{S3 --- 参数更新。} 更新参数子集 (嵌入 / 适配器 / 部分 / 全量 ---
从窄到宽分级), 使 σ 在一个短序列中直接映射到其结果, 同时保持其他能力
(受限范围, 或全量更新后验证)。

\textbf{S4 --- 语义验证。} 在留出集上以完整输出序列比对 (唯一可信指标)
验证快路径输出与慢路径输出一致。只有通过的宏才被安装。

\textbf{S5 --- 运行时路由。} 当 σ 的执行成本达到阈值时路由到快路径;
验证失败回退慢路径。快路径是优化, 不是替代。

\textbf{与蒸馏的关键区别。} 蒸馏中, (可能是外部的)
教师向学生迁移知识。这里教师和学生是\emph{同一模型}的慢路径与快路径:
构造通道监督自己的直觉通道。无新语义机制; 宏从已有能力编译而来。

\subsection{4 实验}\label{ux5b9eux9a8c}

全部经官方 run\_exp 判定路径 (完整序列重建)。模型: exp02\_supervised\_s2
(2 层, dim-64, \textless1 MB, 完整套件判定 1.000)。定位: P1
确立该模型泛化 (宽度/进制/匿名化);
本文确立其\emph{已泛化}的能力可被重新编译为更廉价的路径 ---
以下实验测量可执行性、成本与跳步, 而非泛化本身 (由 P1 覆盖)。

\subsubsection{4.1 零样本语义闭包
(C2)}\label{ux96f6ux6837ux672cux8bedux4e49ux95edux5305-c2}

\begin{longtable}[]{@{}llll@{}}
\toprule\noalign{}
输入 & 构造 & 形式 & 直觉 (判定) \\
\midrule\noalign{}
\endhead
\bottomrule\noalign{}
\endlastfoot
12+34 & ✓ & ✓ & ✓ \\
123456789+987654321 & ✓ & ✓ & ✓ \\
999\ldots9 (20 位)+1 & ✓ & ✓ & ✓ \\
\textbf{999\ldots9 (2000 位)+1} & ✓ & ✓ & \textbf{1.000} \\
\end{longtable}

构造与形式通道在全部四个输入上完全一致 (完整输出序列); 直觉 (模型)
在相同输入上判定准确率 1.000, 包括 2000 位进位链 ---
对该特定算子零训练。语义先于训练: 仅定义链就使 σ 可执行。

\subsubsection{4.2 成本画像 (C1,
C4)}\label{ux6210ux672cux753bux50cf-c1-c4}

\begin{longtable}[]{@{}
  >{\raggedright\arraybackslash}p{(\linewidth - 8\tabcolsep) * \real{0.2000}}
  >{\raggedright\arraybackslash}p{(\linewidth - 8\tabcolsep) * \real{0.2000}}
  >{\raggedright\arraybackslash}p{(\linewidth - 8\tabcolsep) * \real{0.2000}}
  >{\raggedright\arraybackslash}p{(\linewidth - 8\tabcolsep) * \real{0.2000}}
  >{\raggedright\arraybackslash}p{(\linewidth - 8\tabcolsep) * \real{0.2000}}@{}}
\toprule\noalign{}
\begin{minipage}[b]{\linewidth}\raggedright
宽度
\end{minipage} & \begin{minipage}[b]{\linewidth}\raggedright
构造 (慢)
\end{minipage} & \begin{minipage}[b]{\linewidth}\raggedright
形式
\end{minipage} & \begin{minipage}[b]{\linewidth}\raggedright
直觉 (快)
\end{minipage} & \begin{minipage}[b]{\linewidth}\raggedright
相对构造加速
\end{minipage} \\
\midrule\noalign{}
\endhead
\bottomrule\noalign{}
\endlastfoot
2 & 0.0005 s & 0.0012 s & 0.0018 s & --- \\
9 & 0.0031 s & 0.0007 s & 0.0043 s & --- \\
\textbf{2000} & \textbf{0.6206 s} & \textbf{0.0695 s} & \textbf{0.0857
s} & \textbf{7× (直觉) / 9× (形式)} \\
\end{longtable}

在 2000 位时, 快路径比构造便宜 7×。构造随位数线性增长;
快路径是短单次前向。诚实警示: 形式引擎更快;
直觉通道的价值在于权重内的端到端执行 (无外部 oracle) 加保持的泛化。

\subsubsection{4.3 跳步泛化 (C4)}\label{ux8df3ux6b65ux6cdbux5316-c4}

在固定获得 (8 轮) 下, 编译的快路径通过跳过构造步骤\textbf{零衰减}泛化
(exp80): 2000 位宽 (10× 训练宽度)、60 进制 (6×)、100\%
数字匿名化、以及宽度 × 进制联合 --- 全部 1.000。获得与跳步泛化解耦:
快路径难以获得但易于外推。

\subsubsection{4.4 跳步的机器形式
(C4)}\label{ux8df3ux6b65ux7684ux673aux5668ux5f62ux5f0f-c4}

\begin{longtable}[]{@{}lll@{}}
\toprule\noalign{}
路径 & Token & 时间 \\
\midrule\noalign{}
\endhead
\bottomrule\noalign{}
\endlastfoot
构造 (逐位展开) & 82 & 0.410 ms \\
直觉 (原子数字 token) & 5 & 0.186 ms \\
\end{longtable}

20 位加法: 快路径将整个数字视为原子 token (5 个) 对比 82 个 token
的逐位展开 --- 16× token 缩减, 2.2× 加速。\textbf{诚实警示}: 原子 token
在模型词表外 (pad id 0); 测量的是序列长度对前向成本的影响 --- 下界,
而非习得原子 token 语义的主张。若模型真正学会原子 token, 快路径将是 O(1)
决策且优势更大。

\subsection{5 消融}\label{ux6d88ux878d}

\begin{longtable}[]{@{}
  >{\raggedright\arraybackslash}p{(\linewidth - 4\tabcolsep) * \real{0.3333}}
  >{\raggedright\arraybackslash}p{(\linewidth - 4\tabcolsep) * \real{0.3333}}
  >{\raggedright\arraybackslash}p{(\linewidth - 4\tabcolsep) * \real{0.3333}}@{}}
\toprule\noalign{}
\begin{minipage}[b]{\linewidth}\raggedright
消融
\end{minipage} & \begin{minipage}[b]{\linewidth}\raggedright
问题
\end{minipage} & \begin{minipage}[b]{\linewidth}\raggedright
状态
\end{minipage} \\
\midrule\noalign{}
\endhead
\bottomrule\noalign{}
\endlastfoot
样本数 (10/20/50) & 多少实例足以编译 & pending \\
更新范围 (嵌入 / 适配器 / 部分 / 全量) & 成本--退化权衡; 原能力保护 &
pending \\
无验证器 / 无回退 & 无安全网时的错误传播 & pending \\
闭包不变性 (EXP-52) & 快路径可执行集 = 慢路径可执行集 & pending \\
\end{longtable}

三通道一致性实验 (exp50) 已提供核心验证证据: 构造 = 形式完全一致
(2/9/20/2000 位) 且直觉 1.000, 包括 2000 位进位链。闭包不变性 (EXP-52)
检验快路径是否精确执行慢路径可执行的输入 --- 目标 (iii) 的语义闭包保证。

\subsection{6 相关工作}\label{ux76f8ux5173ux5de5ux4f5c}

\begin{itemize}
\tightlist
\item
  \textbf{程序到神经编译} (Learning to Compile Programs to Neural
  Networks, ICML 2024): 从行为数据为\emph{外部}程序构建神经替代;
  需要程序规约与大量数据。区别: 这里的''程序''是模型自身的构造路径;
  教师是模型自身; 无程序分析。相关综述视角见 P1 §9。
\item
  \textbf{知识蒸馏} (Hinton et al., 2015): 教师→学生,
  通常是独立/更大的教师。区别: 同一模型, 两个通道; 无外部教师;
  等价验证带回退。
\item
  \textbf{推理编译 / 捷径蒸馏} (diffusion, 2025):
  少样本捷径蒸馏压缩长的正确路径。表面相似; 区别: 带验证器 +
  回退的语义等价, 无语义闭包扩展。
\item
  \textbf{少样本微调}: 通用; 本工作是''能力已存在于慢路径''的特殊情形,
  少量实例是\emph{编译}而非\emph{学习}。
\item
  \textbf{持续 / 终身学习}:
  在不遗忘旧能力下增加新能力。本工作的更新范围限制 (S3) 与验证 (S4)
  是遗忘缓解机制; 消融 (§5) 测量权衡。
\item
  \textbf{机制性 / 手工编译 transformer} (如手工编译的 WASM 解释器):
  手工指定权重解释语言。我们的是\emph{习得}编译 --- 梯度训练产生宏,
  无手工设置权重。
\item
  \textbf{持续学习 / 灾难性遗忘}: 微调有遗忘风险; 我们的受限更新 + 验证
  + 回退是针对性回应, 更新范围消融 (§5) 显式检验。
\item
  \textbf{长度泛化文献} (位置编码 / 循环工程):
  处理固定表示下的长输入推理; 本工作的快路径在通道层面消除构造长度依赖,
  而非表示层面。
\end{itemize}

\subsection{7 局限}\label{ux5c40ux9650}

\begin{enumerate}
\def\labelenumi{\arabic{enumi}.}
\tightlist
\item
  \textbf{可编译集。} 只有可由已支持运算组合表达的算子可编译 (S1
  前提)。新语义不在范围内: 宏增加速度, 不增加意义。
\item
  \textbf{慢路径正确性。} 编译质量继承构造通道的正确性; 参考求值器 bug
  会传播进监督。验证缓解但不创造正确性。
\item
  \textbf{验证覆盖。} 等价在留出样本上检验, 而非证明;
  完整保证需要编译权重的形式验证 (未来, Lean 基)。
\item
  \textbf{更新范围权衡。} 宽更新有破坏已学能力的风险;
  窄更新可能限制快路径表达力。经验前沿尚未绘制 (消融 pending)。
\end{enumerate}

\subsection{8 结论}\label{ux7ed3ux8bba}

一个已沿构造路径执行复合算子的模型, 拥有快速执行它所需的一切:
少量样本把构造通道编译为同一模型内语义等价的推理快路径 ---
可验证、带回退、不扩展语义闭包。在一个 2 层 \textless1 MB transformer
上: 2000 位进位链所有通道正确, 快路径在宽度 2000 时便宜 7×,
跳步零衰减泛化, 宏的机器形式可见为 82→5
token。神经宏编译是源码级宏的权重级类比: 语义是模型自己的; 只增加速度。

\subsection{附录}\label{ux9644ux5f55}

\begin{itemize}
\tightlist
\item
  A. 完整实验表: docs/paper\_data/results.csv。
\item
  B. 路径识别示例: 加法/根/幂的定义链展开。
\item
  C. 成本协议: 构造 (token\_iterator), 形式 (引擎), 直觉 (模型前向); 每
  1000 样本墙钟时间。
\item
  D. 专利 M1 独立权利要求骨架 (语义路径编译)。
\end{itemize}

\subsection{统一框架内的交叉引用}\label{ux7edfux4e00ux6846ux67b6ux5185ux7684ux4ea4ux53c9ux5f15ux7528}

本文是《表示、逻辑、直觉的统一框架》(纲领:
\texttt{Unified\_Framework\_Representation\_Logic\_Intuition.md})
的一条腿。三条腿相互适配:

\begin{longtable}[]{@{}lll@{}}
\toprule\noalign{}
腿 & 论文 & 主张 \\
\midrule\noalign{}
\endhead
\bottomrule\noalign{}
\endlastfoot
形式 → 构造 & P1 (本文) & 正确语法 → 极快训练; 正确构造 → 极强泛化 \\
构造 → 直觉 & P0 (三通道等价) & 构造 = 形式 = 直觉; 直觉是编译的构造 \\
构造 → 直觉 (方法) & P2 (神经宏编译) & 把慢构造路径编译为快推理路径 \\
\end{longtable}

循环闭合: 形式指导构造 (P1), 构造获得直觉 (P0/P2), 直觉搜索形式
(快路径外推并再表达于统一形式)。统一结论 --- 形式指导逻辑构造,
构造获得闪念直觉, 直觉穷举搜索形式 ---
在纲领与每篇论文的统一结论中陈述。

\subsection{参考文献}\label{ux53c2ux8003ux6587ux732e}

{[}准备中; 完整记录待核实。核心条目:{]}

\begin{itemize}
\tightlist
\item
  Hinton, G., Vinyals, O., Dean, J. 2015. Distilling the Knowledge in a
  Neural Network. arXiv:1503.02531.
\item
  ICML 2024. Learning to Compile Programs to Neural Networks
  (完整引用待核实)。
\item
  P1 姊妹论文: Generalization as Formalized Induction (2026)。
\item
  Diffusion 捷径蒸馏 (2025; 引用待定)。
\item
  Barendregt, H. 1984. The Lambda Calculus (定义链与 λ 基础)。
\end{itemize}

\begin{center}\rule{0.5\linewidth}{0.5pt}\end{center}

\subsection{统一结论 (2026-08-11
增补)}\label{ux7edfux4e00ux7ed3ux8bba-2026-08-11-ux589eux8865}

当我们在统一的形式化注意力语法结构下,
最终不可避免地、必然地通过神经网络清晰、直观、可解释地观察到直觉与构造的边界时,
我们就可以有充足的信心: 用统一的注意力形式语言指导逻辑构造, 用精确的合成
CoT 逻辑构造获得 \emph{闪念直觉},
用蒸馏的直觉上下文穷举注意力形式。形式、构造、直觉在统一框架下相互适配,
组成再不分彼此的教学、训练、推理的反馈迭代管线 --- 这也许是我们通往 ASI
的众多路径中, 相对较快的一条捷径。

\end{document}
